\section{Returning an ordinary Type II point to the original selected component}
\label{sec:cubic-selected-component}

This section proves a pointwise comparison with the source's
\emph{original} local ideal $I_t$.  It does not identify all
tangent spaces on that ideal.  The distinction permits the
original Type II modularity hypothesis to be used for a point
obtained by removing its cubic endpoint character, once the
other original point conditions have been checked.

Fix the source datum $D=\texttt{Loc.Dat3}$ over the integers
$\mathcal O$ of a finite extension $E/\mathbb Q_\ell$,
and its framed universal ring $R$.  Write
\[
 J=I_Y(D),\qquad P=I_t(D),\qquad
 r_0=r_0(\eta,\Psi),\qquad H_L=\ker\eta\cap\ker\tau.
                                                        \tag{G48.1}
\]
All uses of the local source theorems retain their stated
\texttt{LocalPublishedTheorems} hypothesis.  In particular,
they are not consequences of the supplied analytic theorem A.

\begin{lemma}[AC-G48: formal families and an unramified change of basepoint]
\label{AC-G48}
The following conclusions hold for the actual ideals in (G48.1).

First, suppose $\Phi:R\to\mathcal O'[[T]]$ is a continuous
map, with $\mathcal O'/\mathcal O$ finite, whose specialization
at $T=0$ is $r_0$.  Suppose its representation modulo $T^n$,
over $E'[T]/T^n$, is crystalline on $H_L$ for every $n\ge1$.
Then
\[
                         \Phi(P)=0.                    \tag{G48.2}
\]

Second, over a fixed coefficient DVR, let $\delta$ be a
continuous character with
\[
 \delta|_I=\Psi|_I,\qquad \overline\delta=\overline\Psi.
                                                        \tag{G48.3}
\]
Replace $\Psi$ by $\delta$, leaving $\eta,\tau,e,K_v$
unchanged.  Under the canonical identification of the two
framed universal rings, their ideals $I_Y$ are equal and
their selected primes $I_t$ are equal.  In particular the
diagonal point $r_0(\eta,\delta)$ belongs to the original
selected component.

Finally, after a finite coefficient extension, the image of
the original selected ideal is contained in the selected
ideal for the extended datum.  Thus membership proved over
a finite coefficient extension implies membership in the
original ideal.
\end{lemma}
\begin{proof}
\proofstep{The exact source input.}
For the basepoint $x_0:R\to\mathcal O$, the source proves
that $P$ is the unique minimal prime of $\sqrt J$ through
$x_0$, and that there is $y\in R$ with
\[
       x_0(y)\ne0,\qquad yP\subset J.                \tag{G48.4}
\]
The stronger inclusion into $J$, rather than merely
$\sqrt J$, is the conclusion of
\path{Development/AutomorphicArithmetic/AutomorphicVarieties/ModularityArgument/CrystallineOrdinaryIdeals/OrdinaryLocusIdealLeKerOfIsBasePoint.lean},
\texttt{saturatedOrdinaryIdeal\_smooth\_at\_residualRep},
via \texttt{SmoothPushforward.smooth}.
The same theorem gives the asserted uniqueness among all
minimal primes containing $J$ and passing through this
basepoint.  Also
\path{CrystallinePointsSaturatedOrdinaryIdealArtinian.lean},
\texttt{crPts\_IY}, gives exactly \texttt{CrPts} for $J$:
an Artinian crystalline deformation of $r_0$ kills $J$.
Its statement wrapper is used by
\texttt{ideal\_cuts\_crystalline\_locus} in
\path{SaturatedOrdinaryIdealCutsOutCrystallineLocus.lean}.
These are the source conclusions used below.

\proofstep{The formal-family assertion.}
Apply \texttt{CrPts} to the composite of $\Phi$ with
$\mathcal O'[[T]]\to E'[T]/T^n$ and the reduction modulo
$T$.  The latter reduction has nilpotent kernel and its
representation is the specified diagonal basepoint.  Thus
every $a\in J$ maps to zero modulo every $T^n$.
The injection $\mathcal O'[[T]]\hookrightarrow E'[[T]]$
and $\bigcap_{n\ge1}T^nE'[[T]]=0$ give $\Phi(J)=0$.
The specialization of $\Phi$ at $T=0$ is the basepoint map
$x_0$: this follows from uniqueness in the framed universal
property, not merely equality of the residual representations.
Thus the constant coefficient of $\Phi(y)$ is the nonzero image
of $x_0(y)$.  Since $\mathcal O'[[T]]$ is a domain,
(G48.4) now gives $\Phi(a)=0$ for every $a\in P$.
This proves (G48.2).  It does not require a converse
crystalline-component comparison at an arbitrary point.

\proofstep{Why the ordinary ideal is unchanged.}
The framed unrestricted deformation functor depends on
$\eta$ and $\overline\Psi$, so (G48.3) identifies its
universal rings and universal representations.
In the literal definitions in
\path{PadicHodgeTheory/IsInertiaInfiniteCyclotomic.lean},
$H_L$, the choices \texttt{s0} and \texttt{fr}, and
the character $\tau$ do not change.  Every occurrence of
$\Psi$ in \texttt{ordIdeal} is evaluated on
$H_L\cap I$.  The coefficients \texttt{cc} evaluate
$\Psi$ at \texttt{s0}, also an inertia element.
The expressions \texttt{tr} use the universal matrix at
\texttt{fr}; \texttt{dd} uses these coefficients and
the \emph{cyclotomic} value at \texttt{fr}.
It contains no value $\Psi(\texttt{fr})$.
Therefore \texttt{ordIdeal}, \texttt{del}, and their
saturation $J$ are exactly the same for both data.

\proofstep{The diagonal interpolation.}
The ratio $a=\delta/\Psi$ is unramified and has values
congruent to $1$.  Let $\alpha=a(\operatorname{Frob}_\ell)
\in1+\mathfrak m_{\mathcal O}$.  The unramified character
$u_T$ with arithmetic Frobenius value $1+T$ is well-defined
over $\mathcal O[[T]]$: on the pro-$\ell$ quotient of
$\widehat{\mathbb Z}$ it is $(1+T)^z$ for
$z\in\mathbb Z_\ell$, using the binomial series.
It is continuous for the $(\mathfrak m_{\mathcal O},T)$
topology.  The diagonal family
\[
                  r_0(\eta,\Psi u_T)                  \tag{G48.5}
\]
has the original residual representation.  Each of its
Artinian coefficient-field truncations is crystalline on
$H_L$: its diagonal characters there are Tate characters
times unramified representations.  Unramified
coefficient-field representations are crystalline even when
their Frobenius is not semisimple, by crystalline descent
from inertia.  An exact primary locator is Brinon--Conrad,
\emph{CMI Summer School Notes on $p$-adic Hodge Theory},
24 June 2009, Proposition~9.3.1, printed p.~146;
the notes are available at
\url{https://virtualmath1.stanford.edu/~conrad/papers/notes.pdf}.
Applying (G48.2) to (G48.5), then evaluating
$T=\alpha-1$, shows that $r_0(\eta,\delta)$ kills $P$.

Apply the source theorem (G48.4) again with $\delta$ as
the base character.  Its selected prime $P_\delta$ is
the unique minimal prime of the \emph{same} $J$ passing
through $r_0(\eta,\delta)$.  Since $P$ is such a minimal
prime, $P=P_\delta$.  This proves the second assertion.
There is no passage between components at a possibly singular
intersection: uniqueness at the new diagonal basepoint is
the explicit reason the two primes agree.

\proofstep{Finite coefficient extension.}
For completeness the natural map of framed universal rings
can be constructed on finite Artinian test rings.
If $A$ is such a local $\mathcal O'$-algebra with residue
field $k'$ and the prescribed residual matrices have entries
in the original residue field $k$, put
$A_0=\operatorname{res}^{-1}(k)$.  Every matrix entry of a
lift, and its inverse determinant, lies in $A_0$.
This is a finite local $\mathcal O$-algebra with residue
field $k$.  Original universality therefore supplies a
map $R\to A_0\to A$.  Uniqueness, followed by inverse
limits, constructs the canonical map to the extended framed
ring $R'$; equivalently the completed coefficient extension
represents the extended framed functor.

Let $J',P'$ be the ideals for the extended original datum.
The defining generators and saturation show that the image
of $J$ is contained in $J'$.  Hence (G48.4) gives
$y\,\operatorname{image}(P)\subset J'\subset P'$.
The element $y$ does not belong to $P'$, because the extended
basepoint evaluates it nontrivially.  Primality of $P'$ gives
$\operatorname{image}(P)\subset P'$, as claimed.
No assertion that radicals or selected primes commute with
arbitrary base change is needed.
\end{proof}

\begin{proposition}[AC-G48-TypeII: an actual selected-component criterion]
\label{AC-G48-TypeII}
Assume $\ell$ is odd and
\[
              \overline\Psi^{2}|_I\ne1.               \tag{G48.6}
\]
Let $r:G_{\mathbb Q_\ell}\to GG_3(\mathcal O')$ be
a continuous integral representation with multiplier $\eta$
and the original residual representation.  Suppose that its
middle line is identically $1$, its mixed entries are zero,
and its endpoint matrix $B(g)$ has determinant $1$ for every
$g$.  Suppose also that:

\begin{enumerate}
\item $r$ is crystalline on $H_L$;
\item it has the source's ordered ordinary inertia conclusion:
over the coefficient field a complete flag is preserved on
$\ker\eta$, and its inertia characters, in order, are
$\Psi,1,\Psi^{-1}$.
\end{enumerate}

Then the local point defined by $r$ kills the original
selected ideal $I_t(D)$.
Condition~2 is the actual conclusion \texttt{TriIn} of
\path{PadicHodgeTheory/CommutingUpperTriangularDistinct.lean},
\texttt{lemma2b}; it specifies an ordered flag, rather
than only a semisimplified inertia type.
\end{proposition}
\begin{proof}
By the last part of AC-G48 we may work over $\mathcal O'$
throughout, and descend the conclusion at the end.
Write $\mathcal O$ for this DVR and $\varpi$ for a
uniformizer in this proof.

\proofstep{A strict integral ordinary basis.}
Put $J_2=\left(\begin{smallmatrix}0&1\\1&0\end{smallmatrix}\right)$,
$D_2=\operatorname{diag}(-1,1)$, and write
$\eta(g)=(-1)^{\epsilon(g)}$.  The elementary identity
\[
 J_2B^{-t}J_2=D_2BD_2\qquad(B\in\mathrm{SL}_2)
\]
shows that
$V(g)=B(g)D_2^{\epsilon(g)}$ is a linear
two-dimensional representation with determinant $\eta$.
This is the endpoint normalization of AC-G45.
On inertia $\eta=1$, so $V$ and the endpoint block agree.

The first line of the prescribed ordinary flag has inertia
character $\Psi$.  Its projection to the middle line is
zero, because $\Psi|_I\ne1$.  It therefore lies in the
endpoint plane; call it $L$.  This is the unique inertia
subline with character $\Psi$, since the other endpoint
character is $\Psi^{-1}$ and these characters are distinct.
For any $g\in G_{\mathbb Q_\ell}$, normality of inertia
and the fact that $\Psi$ is a character of the full group
show that $V(g)L$ has that same inertia character.  Thus
$L$ is stable under the full group.

The saturated lattice $L\cap\mathcal O^2$ is a direct
summand.  Its nonzero reduction is an inertia eigenline with
character $\overline\Psi$.  By (G48.6) the original
residual endpoint characters are distinct, so that line is
the first coordinate line.  It has a generator $(1,u)$ with
$u\in\varpi\mathcal O$.  The matrix
\[
              N=\begin{pmatrix}1&0\\u&1\end{pmatrix}
               \in\mathrm{SL}_2(\mathcal O)           \tag{G48.7}
\]
is congruent to $1$ and gives an upper triangular basis.
In that basis write
\[
 N^{-1}V(g)N=
 \begin{pmatrix}\eta(g)\delta(g)&b(g)\\
                 0&\delta(g)^{-1}\end{pmatrix}.       \tag{G48.8}
\]
Here $\delta$ is a continuous unit-valued character,
$\delta|_I=\Psi|_I$, and
$\overline\delta=\overline\Psi$ on the full group.
Also $b(g)\in\varpi\mathcal O$ for every $g$.
These assertions follow by reducing the strict basis change
and comparing the diagonal and off-diagonal entries with
the original residual representation.  In particular
$\delta/\Psi$ is an unramified character congruent to $1$.

\proofstep{Contract the crystalline extension.}
For a formal variable $T$, replace $b(g)$ in (G48.8) by
$T b(g)/\varpi$ and call the result $V_T(g)$.
This is a representation over $\mathcal O[[T]]$:
the upper-right entry is the extension cocycle of the two
fixed diagonal characters, and its cocycle equation is
linear.  It is continuous, has determinant $\eta$, and
has the prescribed residual representation.
Convert $V_T$ back to the semidirect endpoint block by
multiplying on the right by $D_2^{\epsilon(g)}$, and
adjoin the constant middle line.  Denote this family by $r_T$.
Its specializations are
\[
 r_T|_{T=0}=r_0(\eta,\delta),\qquad
 r_T|_{T=\varpi}=\operatorname{conj}_{N^{-1}}r.       \tag{G48.9}
\]

The second identity is an actual semidirect conjugation:
for $N\in\mathrm{SL}_2$,
\[
 N^{-1}V(g)N D_2^{\epsilon(g)}
 =N^{-1}B(g)\,\theta_{\eta(g)}(N),
 \qquad \theta_{\eta(g)}(N)=D_2^{\epsilon(g)}ND_2^{\epsilon(g)}.
\]
Thus the linear basis change and the source endpoint
conjugation have precisely the same matrix component.

Each truncation over $E[T]/T^n$ is crystalline on $H_L$.
Indeed, restrict first to $H_L$, where the representation
is linear.  The original two-dimensional extension is
crystalline by hypothesis, and its diagonal characters are
crystalline.  Over $B_{\mathrm{cris}}$ its extension
cocycle is a coboundary.  Multiplying that splitting
cocycle by $T/\varpi$ splits the entire truncated extension
over $B_{\mathrm{cris}}\otimes E[T]/T^n$.
The middle line is already crystalline.  Thus the period
criterion holds for every truncation, including the
nonreduced ones.

By AC-G48, replacing the source base character $\Psi$ by
$\delta$ leaves both $J$ and $P$ unchanged.  Apply the
formal-family assertion to $r_T$, now at this new diagonal
basepoint.  It kills $P$.  The specialization $T=\varpi$
is defined on all formal power series and therefore shows
that $\operatorname{conj}_{N^{-1}}r$ lies on $P$.
Finally strict integral conjugation preserves $P$:
this is exactly
\path{PadicHodgeTheory/RtConj.lean}, \texttt{rtConj},
proved in
\path{ModularityArgument/CrystallineOrdinaryIdeals/ExistsConjLift.lean},
\texttt{GlobalTheory.OrdLoc.rtConj}.
Apply it to the endpoint embedding of $N$, which is
congruent to $1$, to recover the original $r$.
Finite coefficient descent from AC-G48 completes the proof.
\end{proof}

\begin{corollary}[AC-G48-cubic: pointwise removal of the cubic endpoint character]
\label{AC-G48-cubic}
Let $\ell\ne3,269$ satisfy the source \texttt{GD} conditions.
Suppose an integral cubic Type II point has the original
ordinary flag characters on inertia above $\ell$, is
crystalline on $H_L$, and has the block and determinant
properties needed for the endpoint descent in AC-G45.
Remove $\chi^j$ from its endpoint block on $G_K$, retain
the middle line $1$, and use the normalized descent with
$r(c)$ having matrix $I$.
If the resulting point has the original residual
representation, then its restriction at $\ell$ belongs
to the original selected local component $I_t$.

Consequently the selected-component field of the original
\texttt{Dat.IsPt} does not require an additional universal
automorphy transport for these ordinary Type II points.
Once its other point fields and its map to the original
global deformation ring are established, the original
pointwise \texttt{LemW} hypothesis can be applied.
This does not construct the modified cubic deformation ring,
prove its Type I automorphy, or establish rigidity of its
mixed rank-four tangent coefficient.
\end{corollary}
\begin{proof}
The cubic character is unramified at $\ell\ne269$.
Its inverse tensor on the endpoint plane preserves the
ordered ordinary flag and its inertia characters.  It
also preserves crystallinity on $H_L$, since unramified
representations and their duals are crystalline.  The
normalized descent of AC-G45 gives precisely the determinant-one
semidirect endpoint block used in AC-G48-TypeII.

The source \texttt{GD.ine} states
$(\operatorname{comp}_\ell\psi)^2=1$, and on inertia
$\Psi=\kappa^e\operatorname{comp}_\ell\psi$.
Hence $\overline\Psi^{2}|_I=
\overline\kappa^{2e}|_I$.
The mod-$\ell$ cyclotomic character is surjective on
inertia.  The source \texttt{GD.gen}, namely
$\ell-1\nmid2e$, proves (G48.6).
The exact definition is
\path{Development/AutomorphicArithmetic/AutomorphicVarieties/HeightAndVanishing/Counts.lean},
\texttt{GD}, fields \texttt{ine} and \texttt{gen}.
Thus AC-G48-TypeII applies.

If the point is initially expressed over the source's
algebraic closure, its local map from the complete
Noetherian framed ring factors through a finite
coefficient field.  To see this, choose a surjection
$\mathcal O[[X_1,\ldots,X_s]]\twoheadrightarrow R$.
The finitely many images of $X_i$ are algebraic over
$E$ and topologically nilpotent for an \texttt{IsLocPt}.
The images of a finite set of generators of the maximal
ideal have a common norm bound strictly less than $1$;
images of all ring elements have norm at most $1$.  Powers
of that ideal therefore have images tending uniformly to
zero, which proves continuity of the map before evaluating
complete power series.
They generate a finite extension $E'$.  Evaluation of each
convergent power series remains in the complete field $E'$;
the integral image lies in $\mathcal O_{E'}$.
This is also supplied directly by
\path{PadicHodgeTheory/PtFinFinite.lean}, \texttt{Loc.ptFin},
applied to the constant dual-number lift.
It explains why the finite-extension formulation of
AC-G48-TypeII covers the source coefficient-field points.

The last assertions retain the logical direction of
\texttt{LemW} in
\path{AutomorphicCases/AutomorphyUnramifiedCoefficients.lean},
and \texttt{Dat.IsPt/IsExt/IsModular} in
\path{UnitaryThreeGroups/RhoKApply.lean}.
They concern actual points of the original ring, rather
than a claim that its global universal datum has already
been changed to include $\chi$.
\end{proof}

\paragraph{Scope relative to AC-G47.}
The new criterion closes selected-component membership for
the specified \emph{ordinary Type II points} after endpoint
untwisting.  AC-G47's arbitrary-point tangent converse
remains unproved.  In particular the argument does not make
the original mixed coefficient of the cubic point a sum of
the two rational rank-two coefficients to which the old
Kato input applies.  That coefficient remains
$V\otimes\operatorname{Ind}_K^{\mathbb Q}\chi^j$.


\paragraph{Geometric dictionary.}
The character modification tensors a decomposed lisse local
system's endpoint summand with an unramified rank-one local
system.  Finite reductions and the corresponding finite
covers use the already audited dictionary of SGA~1,
Expos\'e~V, \S\S4 and~7.  The ordinary flag is additional
local linear data, and the selected formal component is
additional deformation data; neither is provided merely
by the finite-cover equivalence.  The contraction above
constructs an actual formal family through a specified
crystalline basepoint and uses its unique selected prime.
This paragraph reuses the earlier categorical readings;
it does not attribute crystalline deformation theory or
an automorphy theorem to SGA~1.
